\documentclass[12pt, a4paper]{article}
\usepackage[a4paper, mag=1000, left=1.5cm, right=1.5cm, top=2cm, bottom=2cm, headsep=0.7cm, footskip=1cm]{geometry}
\usepackage[T2A]{fontenc}
\usepackage[utf8]{inputenc}
\usepackage[english,russian]{babel}
\usepackage[colorlinks=true, citecolor=blue, urlcolor=blue, linkcolor=blue]{hyperref}
\usepackage{graphicx}
\usepackage{physics}
\usepackage{amsmath,amssymb}
\usepackage{indentfirst}
\usepackage{float}
\usepackage{multirow}
\usepackage{hhline}
\usepackage{pgfplots}
\pgfplotsset{compat=1.9}
\usepackage{bigstrut}
\usepackage{fancyhdr}

\setlength{\headheight}{14.5pt}
% Укажите ваш город и текущий год в колонтитуле
\fancyfoot[C]{\text{Ярославль, 2026}}

\begin{document}
	\thispagestyle{fancy}
	
	% --- ТИТУЛЬНЫЙ БЛОК ---
	\begin{center}
		\textsc{Ярославский государственный университет им. П.Г. Демидова\\[3mm]
			Физико-технический факультет\\Электроника и наноэлектроника} \\[3mm]
		% Убедитесь, что файл логотипа (например, logoyargu.png) лежит в одной папке с проектом
		\includegraphics[scale=0.08]{logoyargu.png}
	\end{center}
	\line(1,0){\textwidth}
	\begin{center}
		\begin{tabular}{rlrl}
			\textbf{Группа:} & \underline{\text{ЭН-21 БО \qquad \qquad \qquad \qquad}} & \textbf{К работе допущен} & \underline{\qquad \qquad \qquad \qquad} \\
			\textbf{Студент:} & \underline{\text{Фамилия И.О.} \qquad \qquad \quad} & \textbf{Работа выполнена} & \underline{\qquad \qquad \qquad \qquad} \\
			\textbf{Преподаватель:} & \underline{\text{Преподаватель}\quad \;\;\;\;\qquad} & \textbf{Отчёт принят} & \underline{\qquad \qquad \qquad \qquad} \\
		\end{tabular}
	\end{center}
	\line(1,0){\textwidth}
	
	\vspace{50mm}
	
	\begin{center}
		\textbf{ОТЧЁТ ПО ЛАБОРАТОРНОЙ РАБОТЕ № <Номер>\\}
		\textbf{«<НАЗВАНИЕ ЛАБОРАТОРНОЙ РАБОТЫ>»}
	\end{center}
	
	\newpage
	% --- КОНЕЦ ТИТУЛЬНОГО БЛОКА ---
	
	\section{Цель работы.}
	<Опишите главную цель работы, например: Исследование... / Измерение...>
	
	\section{Задачи, решаемые при выполнении работы.}
	1) <Задача 1>
	
	2) <Задача 2>
	
	3) <Задача 3>
	
	\section{Объект исследования.}
	<Что именно исследуется в работе>
	
	\section{Метод экспериментального исследования.}
	<Краткое описание метода, используемых подходов или принципов>
	
	\section{Рабочие формулы и исходные данные.}
	
	<Описание рабочей величины>:
	\begin{equation}
		V = \frac{A}{B},
		\label{eq:formula1}
	\end{equation}
	где \(A\) — <описание параметра>, \(B\) — <описание параметра>.
	
	Относительная погрешность:
	\begin{equation}
		\delta = \frac{|V_{\text{эксп}} - V_{\text{табл}}|}{V_{\text{табл}}} \cdot 100\%.
		\label{eq:delta}
	\end{equation}
	
	Справочные/табличные значения: <значения, если применимо>.
	
	\section{Измерительные приборы.}
	\begin{table}[H]
		\centering
		Таблица 1 - Измерительные приборы
		\begin{tabular}{|c|c|c|c|c|} \hline
			№ & Наименование & Тип & Диапазон & Погрешность \\ \hline
			1 & <Прибор 1> & <Тип> & <Диапазон> & \(\pm <\text{Погрешность}>\) \\ \hline
			2 & <Прибор 2> & <Тип> & <Диапазон> & \(\pm <\text{Погрешность}>\) \\ \hline
		\end{tabular}
	\end{table}
	
	\section{Схема установки.}
	\begin{figure}[H]
		\begin{center}
			\begin{minipage}{0.45\textwidth}
				\centering
				% Замените '1' на имя вашего файла с картинкой схемы установки
				%\includegraphics[height=4cm]{1}
				\mbox{Рисунок 1 - Схема экспериментальной установки}
			\end{minipage}
		\end{center}
	\end{figure}
	
	\section{Ход работы.}
	
	\subsection{Результаты прямых измерений.}
	
	\begin{table}[H]
		\centering
		\begin{tabular}{|c|c|c|c|}
			\hline
			№ & \(X\), <ед.> & \(Y\), <ед.> & \(Z\), <ед.> \\
			\hline
			1 & ... & ... & ... \\
			2 & ... & ... & ... \\
			3 & ... & ... & ... \\
			\hline
		\end{tabular}
		\caption{Результаты измерений}
		\label{tab:results}
	\end{table}
	
	\subsection{Расчёты.}
	
	По формуле \eqref{eq:formula1}:
	\[
	V_1 = \dots, \quad V_2 = \dots
	\]
	
	Среднее значение:
	\[
	V_{\text{ср}} = \dots
	\]
	
	\subsection{График зависимости.}
	
	\begin{figure}[H]
		\centering
		\begin{tikzpicture}
			\begin{axis}[
				width=0.8\textwidth,
				height=0.45\textwidth,
				xlabel={\(X\), <ед.>},
				ylabel={\(Y\), <ед.>},
				grid=both,
				legend pos=north west,
				]
				\addplot[blue, thick, mark=*] coordinates {
					(1, 2) (2, 4) (3, 6) (4, 8)
				};
				\addlegendentry{\(Y = f(X)\)}
			\end{axis}
		\end{tikzpicture}
		\caption{График исследуемой зависимости}
		\label{fig:graph}
	\end{figure}
	
	\section{Окончательные результаты.}
	
	Полученное значение:
	\[
	V_{\text{ср}} = <\text{Значение}> \pm <\text{Погрешность}>.
	\]
	
	Относительная погрешность:
	\[
	\delta \approx <\text{Значение}>\%.
	\]
	
	\section{Выводы.}
	
	В ходе работы <краткий итог проделанной работы>. Полученные результаты <согласуются / не согласуются> с теоретическими/табличными значениями. Основные источники погрешностей: <факторы>. Цель работы достигнута.
	
\end{document}